%% CV mẫu tiếng Việt - Phạm Thị Phương (팜 티 프엉)
%% Dùng khi tin tuyển dụng viết bằng tiếng Việt (xem 10-vietnam-market.md).
%% Bố cục giống main_example.tex; chỉ khác tiêu đề mục, thứ tự họ tên và
%% định dạng ngày theo thói quen nhà tuyển dụng Việt Nam.
%% Latin Modern (font mặc định của moderncv khi chạy lualatex) có đủ dấu tiếng Việt.
%%
%% SETUP: chạy /setup để điền thông tin thật.
%% Compile with: lualatex main_example_vi.tex
%% (pdflatex often fails on modern MiKTeX with fontawesome5 font-expansion errors.)

\documentclass[11pt,a4paper,sans]{moderncv}
\moderncvstyle{banking}
\moderncvcolor{blue}

% Force the name and section headings to render in moderncv blue (color1).
% Default banking leaves them black: moderncvstylebanking.sty's \colorlet
% copies (not aliases) the pre-scheme accent colour, so the name colours are
% frozen before \moderncvcolor runs. Re-let them after. \namefont is the hook
% every name-style macro routes through, so this also works on moderncv 2.3.1
% (Debian/Ubuntu apt), which has no \firstnamestyle/\lastnamestyle at all.
\renewcommand*{\namefont}{\fontsize{34}{36}\bfseries\upshape}
\colorlet{firstnamecolor}{color1}
\colorlet{lastnamecolor}{color1}
\colorlet{namecolor}{color1}
\renewcommand*{\sectionstyle}[1]{{\sectionfont\color{color1}#1}}

\usepackage[utf8]{inputenc}
% pdflatex fallback only (the documented engine is lualatex, which skips this
% branch). Without T1 font encoding pdflatex builds accented letters with
% \accent, and the PDF text layer stores them decomposed - `e` + U+0300 rather
% than U+00E8 - so an ATS keyword match on "Genève" fails while the page looks
% right. moderncv 2.5 loads T1 itself under pdflatex; 2.3.1 (Debian/Ubuntu apt)
% does not. \ifpdftex comes from iftex, which every moderncv version loads.
\ifpdftex\usepackage[T1]{fontenc}\fi
% moderncv loads hyperref itself in an \AtEndPreamble hook, so \hypersetup
% must go in an \AtEndPreamble of our own: on moderncv < 2.4 a top-level
% \usepackage{hyperref} clashes with the class's own
% \RequirePackage[unicode]{hyperref}. From 2.4.0 the class passes its options
% through \PassOptionsToPackage instead, which is what removes that clash.
\AtEndPreamble{\hypersetup{
    colorlinks=true,
    linkcolor=blue,
    filecolor=magenta,
    urlcolor=blue,
    pdftitle={Phạm Thị Phương - CV},
    % Keep pdfpagemode=UseNone: this block runs after moderncv's own
    % \AtEndPreamble (moderncv.cls sets pdfpagemode there), so a FullScreen
    % value here would win and open every CV in fullscreen presentation mode.
    pdfpagemode=UseNone,
}}
\usepackage[scale=0.80]{geometry}
\usepackage{import}

% personal data
% Thứ tự tiếng Việt: \name{<Họ và tên đệm>}{<Tên>}, ví dụ \name{Nguyễn Văn}{An}
\name{Phạm Thị}{Phương}
% If you have no address to list, DELETE this whole line. \address{}{}{} fails
% with "There's no line here to end" on every moderncv version.
\address{Phú Đông, Nhơn Trạch, Đồng Nai}{}{}
\phone[mobile]{0963 569 816}
\email{Phamphuong121290@gmail.com}

% Ảnh chân dung (tuỳ chọn - phổ biến ở CV Việt Nam, nhưng công ty nước ngoài
% thường không yêu cầu). Bỏ dấu % và đặt file ảnh cạnh file .tex:
% \photo[64pt][0.4pt]{photo}

\begin{document}

\makecvtitle

% ============================================================
%     MỤC TIÊU NGHỀ NGHIỆP / GIỚI THIỆU
% ============================================================

\vspace{6pt}
\small{Kế toán Tổng hợp (총괄 회계 담당자) với hơn 11 năm kinh nghiệm thực chiến trong công tác kế toán tài chính, kê khai \& hoàn thuế GTGT, lập báo cáo tài chính và kiểm toán nội bộ tại doanh nghiệp sản xuất xuất khẩu và doanh nghiệp FDI Hàn Quốc. Tốt nghiệp song bằng Cử nhân Kế toán (ĐH Kinh tế Quốc dân) và Cử nhân Hàn Quốc học (ĐH KHXH\&NV TP.HCM), sử dụng thành thạo tiếng Hàn chuyên ngành cùng các phần mềm Ecount ERP, Smart Pro và Fast.}

% ============================================================
%     NĂNG LỰC CHÍNH
% ============================================================

\section{Năng lực chính}
\vspace{1pt}
\begin{itemize}

\item \textbf{Kế toán Tổng hợp \& Báo cáo Tài chính}: Hạch toán tổng hợp toàn phần hành, lập Báo cáo tài chính tháng/quý/năm, lập báo cáo phân tích quản trị và thực hiện kiểm toán nội bộ định kỳ.

\item \textbf{Kế toán Thuế \& Pháp lý Doanh nghiệp}: Kê khai thuế GTGT, TNCN, TNDN; lập bộ hồ sơ hoàn thuế GTGT hàng quý; trực tiếp tiếp đoàn và giải trình số liệu với Cơ quan Thuế, Quản lý thị trường, VSATTP và PCCC.

\item \textbf{Quản lý Dòng tiền, Công nợ \& Kiểm soát Chi phí}: Quản lý quỹ tiền mặt và tiền gửi ngân hàng, lập báo cáo thu chi hàng ngày cho Ban Giám đốc; đối chiếu công nợ phải thu/phải trả (AR/AP); kiểm tra báo giá, PO, hợp đồng mua bán và hóa đơn điện tử.

\item \textbf{Tiền lương, BHXH \& Nhân sự Nước ngoài}: Tính lương hàng tháng, trích nộp BHXH bắt buộc, quyết toán thuế TNCN, phụ trách tuyển dụng nhân sự và thủ tục xin Giấy phép lao động (Work Permit) cho chuyên gia nước ngoài.

\item \textbf{Phần mềm Kế toán ERP \& Ngoại ngữ}: Thành thạo Ecount ERP, Smart Pro, Fast Accounting, Microsoft Excel; giao tiếp và làm việc trực tiếp bằng Tiếng Hàn (Cử nhân Hàn Quốc học) và Tiếng Anh (trình độ B2).

\end{itemize}

% ============================================================
%     KINH NGHIỆM LÀM VIỆC
% ============================================================

\section{Kinh nghiệm làm việc}
\vspace{3pt}
\begin{itemize}

% --- Vị trí gần nhất ---
\item{\cventry{07/2023 - nay}{Kế toán Tổng hợp (총괄 회계 담당자)}{Fu Food \& Woori Food, Fiktech}{TP. Hồ Chí Minh / Đồng Nai}{}{\vspace{1pt}
\begin{itemize}
    \item {Quản lý toàn diện thu chi tiền mặt và tiền gửi ngân hàng, theo dõi công nợ phải thu/phải trả (AR/AP) và xuất hóa đơn GTGT đầu ra.}
    \item {Kiểm tra tờ khai thuế hàng tháng/quý và lập Báo cáo tài chính; làm việc trực tiếp với Cơ quan Thuế, Quản lý thị trường, VSATTP và PCCC.}
    \item {Phụ trách công tác tuyển dụng, nhân sự, tính lương hàng tháng, quản lý Bảo hiểm xã hội và thực hiện thủ tục xin Giấy phép lao động (Work Permit) cho người nước ngoài.}
\end{itemize}}}

\vspace{3pt}

% --- Vị trí trước ---
\item{\cventry{01/2015 - 06/2023}{Nhân viên Kế toán (회계원)}{Công ty TNHH Giày Tuấn Việt (Giay Tuan Viet Co., Ltd.)}{Đồng Nai}{}{\vspace{1pt}
\begin{itemize}
    \item {Kiểm tra báo giá, đơn đặt hàng (PO), hợp đồng mua bán; kiểm tra tính hợp lệ của hóa đơn chứng từ phục vụ kê khai thuế GTGT hàng tháng và lập bộ hồ sơ hoàn thuế GTGT hàng quý.}
    \item {Lập và nộp Báo cáo tài chính hàng quý, lập báo cáo phân tích quản trị và thực hiện kiểm toán nội bộ định kỳ.}
    \item {Quản lý quỹ tiền mặt, lập báo cáo thu chi hàng ngày gửi Ban Giám đốc; theo dõi và đối chiếu công nợ phải thu/phải trả (AR/AP).}
    \item {Lập bảng lương và BHXH hàng tháng; tính thuế Thu nhập cá nhân (TNCN) và nộp tờ khai thuế đúng quy định.}
\end{itemize}}}

\end{itemize}

% ============================================================
%     HỌC VẤN
% ============================================================

\section{Học vấn}
\vspace{1pt}
\begin{itemize}

\item{\cventry{2016-2018}{Cử nhân ngành Hàn Quốc học (한국학과)}{Trường Đại học Khoa học Xã hội và Nhân văn -- ĐHQG TP.HCM}{TP. Hồ Chí Minh}{}{\vspace{1pt}
Xếp loại tốt nghiệp: Khá (Good). Đào tạo chuyên sâu ngôn ngữ Hàn Quốc thương mại và văn hóa doanh nghiệp Hàn Quốc.
}}

\vspace{3pt}

\item{\cventry{2011-2013}{Cử nhân ngành Kế toán (회계학과)}{Trường Đại học Kinh tế Quốc dân}{Hà Nội}{}{\vspace{1pt}
Xếp loại tốt nghiệp: Khá (Good). Chuyên ngành Kế toán doanh nghiệp, Kế toán thuế và Kiểm toán.
}}

\end{itemize}

% ============================================================
%     CHỨNG CHỈ (tuỳ chọn)
% ============================================================

\section{Chứng chỉ}
\vspace{1pt}
\begin{itemize}
\item{\textbf{Chứng chỉ Tin học Ứng dụng (컴퓨터 활용 능력 자격증)} -- Hoàn thành 01/2019.}
\item{\textbf{Chứng chỉ Tiếng Anh trình độ B2 (영어 B2 자격증)} -- Hoàn thành năm 2018.}
\end{itemize}

% ============================================================
%     NGOẠI NGỮ
% ============================================================

\section{Ngoại ngữ}
\vspace{1pt}
\begin{itemize}
\item {Tiếng Việt (bản ngữ), Tiếng Hàn (Thành thạo công việc -- Cử nhân Hàn Quốc học), Tiếng Anh (Chứng chỉ B2 -- 2018).}
\end{itemize}

% ============================================================
%     NGƯỜI THAM CHIẾU
% ============================================================

\section{Người tham chiếu}
\vspace{1pt}
\begin{itemize}
\item{Sẵn sàng cung cấp khi được yêu cầu.}
\end{itemize}

\end{document}
